% paper/sections/09b_split_ppe.tex
%
% §9.3  分相 PPE 解法：密度比非依存な高精度 PPE 戦略
% ═══════════════════════════════════════════════════════════════════════════════
\subsection{分相 PPE 解法：密度比非依存な高精度 PPE 戦略}
\label{sec:split_ppe_framework}
% backward-compat labels
\label{sec:split_ppe_motivation}%
\label{sec:gfm_numerical_motivation}%
% ═══════════════════════════════════════════════════════════════════════════════

\noindent\textbf{本節の位置づけ：}
§~\ref{sec:ccd_poisson_matrix} で構築した CCD 離散化は
等密度ポアソン方程式に対して $\Ord{h^6}$ の設計精度を達成する．
しかし二相流の PPE は変密度演算子 $\bnabla\cdot(\rho^{-1}\bnabla p) = q$ であり，
界面を横切るステンシルが $\Ord{1}$ の打ち切り誤差を生じる．
本節では，この本質的制約を回避する\textbf{分相 PPE 解法}を本稿の主要 PPE 戦略として定式化する．
一言でいえば，全領域を一本の変密度 Poisson 問題として解くのをやめ，
液相・気相では定数密度 Poisson 問題を解き，界面では圧力ジャンプ条件で両相を縫い合わせる，
という設計である．

\noindent\textbf{本節で選ぶもの：}
分相 PPE は単独の名前ではなく，次の 5 つを同時に満たす pressure stack を指す．
\[
  \text{相内定数密度 Poisson}
  \;+\;
  j_{gl}=p_g-p_l=-\sigma\kappa_{lg}
  \;+\;
  G_\Gamma(p;j)=G(p)-B(j)
  \;+\;
  \mathrm{HFE}
  \;+\;
  \text{大域ゲージ}
\]
したがって，本節で退ける対象は
「一括 smoothed-Heaviside PPE を高密度比 production とみなすこと」，
「CSF 体積力を最終的な表面張力閉包とみなすこと」，
「regular pressure field $J$ を後で足す \texttt{jump\_decomposition} を
affine jump と同等視すること」である．

% ─────────────────────────────────────────────────────────────────────────────
\subsubsection{変密度一括解法の本質的制約}
\label{sec:monolithic_ppe_limitation}
% ─────────────────────────────────────────────────────────────────────────────

Smoothed Heaviside による一括解法では，遷移層（幅 $\varepsilon = 1.5h$）内で
$1/\rho$ が $\rho_l/\rho_g$ 倍のジャンプを持つ．
界面を横切る CCD/FD ステンシルは $\Ord{1}$ の打ち切り誤差を生じ，
PPE のグリーン関数の非局所性により解全体に伝播する
（§~\ref{sec:ppe_condition_number}）．
この誤差は反復回数 $k$ の増大や緩和係数 $\omega$ では回復不能であり，
smoothed Heaviside 一括解法の\textbf{本質的制約}である：
%
\begin{itemize}[noitemsep]
  \item 中程度密度比：高解像度側で格子収束がプラトー化
  \item 中-高密度比：頑健な DC 発散閾値は仮定せず，界面型密度場で設計精度は回復しない
  \item 水--空気級高密度比：一括 smoothed Heaviside のみでの高精度計算は未保証
\end{itemize}
%
したがって本章では，一括 smoothed-Heaviside PPE の高密度比利用を理論保証外とする．
重要な点として，中程度の密度比でも界面近傍の精度は
CCD の設計精度 $\Ord{h^6}$ から\textbf{大幅に劣化}しており，
低密度比だから一括解法で十分ということはない．

\noindent\textbf{条件数の定量的見積もり：}
smoothed Heaviside 一括演算子 $\bnabla\cdot(\rho^{-1}\bnabla\cdot)$ の条件数は
$\kappa_2 \approx (\rho_l/\rho_g)\cdot h^{-2}$ のオーダーで成長する．
水--空気級の高密度比（$\rho_l/\rho_g \sim 10^3$），$h = 1/64$ では
$\kappa_2 \sim 10^3\cdot 4096 \approx 4\times 10^{6}$ となり，
ゲージ射影後の非零固有空間でも CG/DC 反復の収束が遅い
（Neumann 零モードそのものは §~\ref{sec:ppe_phase_gauge_pin} で除去する）．
分相 PPE では各相内 $\kappa_2 \approx h^{-2}\sim 10^{3}$ に縮退する
（高密度比 PPE の解析と緩和戦略は~\cite{Guermond2006,GuermondQuartapelle1998}；
界面ジャンプを陽的に扱う Immersed Interface Method の系統は~\cite{LeVequeLi1994,Li2022}）．

% ─────────────────────────────────────────────────────────────────────────────
\subsubsection{分相 PPE 解法の定式化}
\label{sec:split_ppe_formulation}
% ─────────────────────────────────────────────────────────────────────────────

分相 PPE 解法の核心は「各相で定数密度のラプラシアンを解く」ことにある．
Smoothed Heaviside 一括解法が $\bnabla\cdot(\rho^{-1}\bnabla p) = q$ を
変動係数 $\rho(\psi)$ のまま全領域で解くのに対し，
分相 PPE 解法は各相 $k\in\{\ell,g\}$ で
%
\begin{equation}
  \nabla^2 p_k = \rho_k \, q \qquad \text{（$\rho_k$ は定数）}
  \label{eq:split_ppe}
\end{equation}
%
を解く（ここで $q = (1/\Delta t)\,\bnabla\cdot\bu^\ast$ は PPE の
Predictor 速度由来の発散ソース；本稿 7-step フローの Step 6 と一致）．
ただし，これは単に各相を切り離して解くという意味ではない．
界面 $\Gamma$ 上では圧力ジャンプと法線フラックス条件を同時に課して，
二つの相の楕円型問題を閉じる：
\begin{align}
  [p]_\Gamma
    &= J_p^\chi
     \quad \text{（未知量 $\chi$ に対応する圧力ジャンプ）},
     \label{eq:split_ppe_pressure_jump} \\
  \left[\frac{1}{\rho}\pd{p}{n}\right]_\Gamma
    &= J_n .
     \label{eq:split_ppe_flux_jump}
\end{align}
\noindent\textbf{全圧力と圧力増分の切り分け：}
以後の式では簡潔のため未知量を $p$ と書くが，
PPE に渡すジャンプは\textbf{解く未知量に対応するジャンプ}でなければならない．
すなわち
\begin{equation}
  J_p^\chi =
  \begin{cases}
    j_{gl}^{\,n+1}, & \chi = p^{n+1}\ \text{を直接解く場合},\\[2mm]
    \delta j_{gl}=j_{gl}^{\,n+1}-j_{gl}^{\,n},
      & \chi=\delta p=p^{n+1}-p^n\ \text{を解く場合}.
  \end{cases}
  \label{eq:split_ppe_unknown_jump_contract}
\end{equation}
この区別を失うと，Young--Laplace 圧力を二重計上するか，
逆に前時刻圧力の自由度に吸収して消してしまう．
本稿の \texttt{affine\_jump} 記法で $j_{gl}$ と書く箇所は，
その step の projection potential が保持すべき
$J_p^\chi$ を指す（全圧力 solve なら $j_{gl}^{n+1}$，
圧力増分 solve なら $\delta j_{gl}$）．
同様に法線速度の連続性から
$J_n=\Delta t^{-1}[\bu^*\cdot\hat{\bm n}_{lg}]$ を用いる．
予測速度が界面法線方向に連続となる production 離散化では $J_n=0$ に退化する．

\noindent\textbf{$J_n$ の物理的正当化：}
連続 NS では界面で法線速度が連続（$[\bu\cdot\hat{\bm n}_{lg}]_\Gamma=0$；非相変化）．
Predictor 離散化が保存的ならば $\bu^*$ も界面で連続となり $J_n\equiv 0$．
非保存的離散化（例：upwind 面補間が界面で非対称）では
$[\bu^*\cdot\hat{\bm n}_{lg}]\sim \Ord{h^p}$ の残差が生じ，
$J_n$ は PPE でこの残差を吸収して corrector 後の $\bu^{n+1}$ を
連続に戻す役割を果たす（$\Ord{h^p}$ 一致時に $J_n=0$ が自動成立）．
HFE（§~\ref{sec:field_extension}）は，これらのジャンプ条件で補正した
片側データを界面越しに延長し，CCD ステンシルが不連続値に触れないようにするための
補助技術である．
$\rho_k$ が定数であるため，各相の離散化演算子は密度ジャンプの影響を一切受けない．

\noindent\textbf{欠陥補正法との関係：}
式~\eqref{eq:split_ppe} は等密度ポアソン方程式であるから，
§~\ref{sec:defect_correction_main} の欠陥補正法（DC $k = 3$）を各相内で適用できる．
等密度条件の下では DC の CCD 高次演算子 $L_H$ と FD 低次演算子 $L_L$ の
スペクトル不整合が最小化されるため，
§~\ref{sec:dc_accuracy} の理論が示す $\Ord{h^7}$（Dirichlet 超収束）
ないし $\Ord{h^5}$（Neumann BC 律速）が\textbf{密度比によらず}回復する．

\begin{tcolorbox}[resultbox, title={分相 PPE + DC $k{=}3$ の相内精度と界面律速}]
\begin{enumerate}[noitemsep]
  \item 各相で定数密度の PPE を独立に解くため，
        界面条件~\eqref{eq:split_ppe_pressure_jump}--\eqref{eq:split_ppe_flux_jump}を満たす
        片側問題として閉じる
  \item DC $k{=}3$ は各相内で $\Ord{h^7}$（Dirichlet）/
        $\Ord{h^5}$（Neumann BC 律速）を達成する
  \item 各相内の離散誤差は密度比に直接依存しない．
        ただしこの主張は\textbf{相内部の Poisson 解}に対するものであり，
        cut-face の閉包まで含めた全体精度を一律に $\Ord{h^6}$ と主張するものではない
  \item 速度補正の勾配 $\bnabla(\delta p)$ には CCD $\Ord{h^6}$（均一格子）
        または $\mathcal{G}^\text{adj}$（非一様格子 + 壁面 BC；§~\ref{sec:fvm_ccd_corrector}）を適用する．
        \texttt{affine\_jump} では同じ面ジャンプ $B_f$ も補正側へ渡し，
        $G_\Gamma=G-B$ として Balanced--Force 整合性を保持する
  \item 界面近傍の最終律速は，FCCD 面ジェット
        （$u_f,u'_f$ は $\Ord{H^4}$，$u''_f$ は $\Ord{H^2}$；
        §~\ref{sec:face_jet_def}），HFE/GFM 延長，曲率評価の最小次数で決まる．
        したがって本節の高次性は
        \emph{相内 $\Ord{h^6}$ / DC 高次解 + 界面閉包の別律速}
        と読むべきである
\end{enumerate}
\end{tcolorbox}

\noindent\textbf{HFE の役割：}
分相 PPE 解法において HFE（§~\ref{sec:field_extension}）は不可欠な基盤技術である．
IPC の $\bnabla p^n$（Predictor ステップ）および
$\bnabla(\delta p)$（Corrector ステップ）が界面を横切る前に，
圧力場をジャンプ補正済み HFE Hermite 補間で滑らかに延長することで，
CCD ステンシルが不連続値に触れることを防止する．

% ─────────────────────────────────────────────────────────────────────────────
\subsubsection{実用 production と研究 DNS 版の切り分け}
\label{sec:pure_fccd_split_ppe}
% ─────────────────────────────────────────────────────────────────────────────

本稿で \textbf{production pressure stack} と呼ぶものは，
\texttt{phase\_separated}，\texttt{affine\_jump}，HFE，DC，
同一 face 契約，大域ゲージを同時に満たす実用閉包である．
これは「全ての部品を純 FCCD だけで閉じる」ことと同義ではない．
DC の base solve，壁面閉包，face-flux projection には低次または FVM 的な補助演算子が残り得るが，
PPE と corrector が共有する face flux 契約を破らない限り production 閉包に含める．

\noindent\textbf{研究 DNS 版：}
純 FCCD 分相 PPE は FVM の制御体積保存を離散不変量として放棄し，
すべての空間演算子を FCCD / HFE / GFM の\textbf{単一高次言語}に統一する
DNS アーキテクチャを採る．質量保存は
$\Ord{\Delta x^4}$--$\Ord{\Delta x^6}$ の漸近的な高次性として扱われ，
制御体積恒等式としては保証されない．
したがって純 FCCD DNS 版は，production 実装の必要条件ではなく，
本節の面演算子と誤差律速を説明する高次極限である．

\noindent\textbf{設計選択の根拠：}
FVM 制御体積保存を放棄する代償として，
(i) CLS 質量保存（§\ref{sec:cls_stages} Stage D）が\textbf{体積質量保存}
（全領域 $\int\psi\,\mathrm{d}V$）を $\Ord{\Delta x^4}$ で独立に担保する，
(ii) 速度場の離散発散フリー性は CCD 発散 $\bnabla\cdot\bu^*$ から構成される
PPE 右辺の射影で保証される，の二経路が保存性を補完する．
FVM では両者が同一の面フラックス恒等式から派生するのに対し，
本稿は「演算子の高次性」と「幾何学的発散フリー性」を独立系統で扱う．
数値検証では，質量保存と BF 整合性を分離して評価する．

\noindent\textbf{アーキテクチャの階層}：
\begin{enumerate}[leftmargin=2em,noitemsep]
  \item 界面幾何：Ridge--Eikonal（\S\ref{sec:ridge_eikonal}）
        で位相・メトリックを分離．
  \item 移流：FCCD + HFE（\S\ref{sec:fccd_advection}）で $\psi$ と
        運動量を共通面ジェット $\mathcal{J}_f$ 上で更新．
  \item 粘性：3 層アーキテクチャ（\S\ref{sec:viscous_3layer}）．
        $\mu\nabla^2\mathbf{u}$ は\textbf{界面越えでは使わない}．
  \item 圧力：分相 FCCD PPE + GFM ジャンプ行（本節）．
  \item 剛性処理：デフェクト補正
        （\S\ref{sec:defect_correction_main}）の高次外側 / 低次内側分解．
\end{enumerate}

% ─────────────────────────────────────────────────────────────────────────────
\subsubsection{位相分離 FCCD PPE：面ジェット随伴ペア}
\label{sec:phase_separated_fccd_adjoint}
% ─────────────────────────────────────────────────────────────────────────────

各相 $q\in\{L,G\}$ で\textbf{FCCD 面ジェット}
$\mathcal{J}_f(p_q)=(p_{q,f},p'_{q,f},p''_{q,f})$ を基本量とする．
面随伴ペア
\begin{equation}
  G_h^\FCCD: \{p_{q,i}\}_\text{node} \to \{p'_{q,f}\}_\text{face},
  \qquad
  D_h^\FCCD = -(G_h^\FCCD)^\ast
  \label{eq:split_ppe_adjoint_pair}
\end{equation}
により相 $q$ の PPE ブロックを SPD 符号規約
$A_h^q\equiv -D_h^\FCCD(\rho_q^{-1}G_h^\FCCD)$ で
\begin{equation}
  A_h^q p_q
    \equiv -D_h^\FCCD\Bigl(\rho_q^{-1}\,G_h^\FCCD\,p_q\Bigr)
    = -\frac{1}{\Delta t}\,D_h^\FCCD\,\mathbf{u}_q^*
  \qquad\text{in }\Omega_q
  \label{eq:split_ppe_fccd_block}
\end{equation}
として組み立てる．$\rho_q$ が各相で\textbf{定数}であるため
$-D_h^\FCCD\rho_q^{-1} G_h^\FCCD$ は
純 Poisson 演算子の FCCD 版となり，
§\ref{sec:ccd_poisson_matrix} の行列構造がそのまま適用される．

\noindent\textbf{eq:split\_ppe との同値性記法}：
本式~\eqref{eq:split_ppe_fccd_block} は~\eqref{eq:split_ppe} の
$\nabla^2 p_q = \rho_q\,q$ と等価：
\textbf{前提}（明示）— $\rho_q$ は相 $q$ 全体で\textbf{格子点に依存しない定数}
（液相・気相それぞれ単一値）；位相マスク境界 $\partial\Omega_q$ は
GFM ゴーストジェット（\S\ref{sec:gfm_ghost_jet_stitch}）で扱い，
ステンシル内に $\rho$ ジャンプは現れない．
従って product-rule 第 2 項
$(\bnabla\rho_q^{-1})\!\cdot G_h p_q$ は消滅し，
$D_h(\rho_q^{-1}G_h p_q) = \rho_q^{-1}D_h G_h p_q = \rho_q^{-1}\nabla_h^2 p_q$
で書き換えると，SPD 形 $-\nabla_h^2 p_q=-\rho_q\,q$ となる．
これは式~\eqref{eq:split_ppe} の Poisson 形 $\nabla^2p_q=\rho_q q$ と同値であり，
符号差は解く線形系を SPD にするための左辺・右辺同時反転にすぎない．

\noindent\textbf{離散随伴性と SPD}：
P-2（§\ref{sec:bf_p2_adjoint}）により
\begin{equation}
  \langle p_q,\,A_h^q p_q\rangle_h
    = \rho_q^{-1}\langle G_h^\FCCD p_q,\,G_h^\FCCD p_q\rangle_h \ge 0,
  \label{eq:split_ppe_spd}
\end{equation}
等号は $G_h^\FCCD p_q \equiv 0$ のみ（相 $q$ 内で一定）．
各相ブロックは定数差モード除去後に SPD で CG/PCG が使える．

% ─────────────────────────────────────────────────────────────────────────────
\subsubsection{同一面係数による投影閉包と浮力残差}
\label{sec:phase_separated_projection_closure}
% ─────────────────────────────────────────────────────────────────────────────

分相 PPE の安定性は，PPE 行列だけでなく，Predictor 残差と Corrector が
同じ面係数を用いることに依存する．これは射影法の
``PPE と速度補正が同じ離散作用素の随伴対でなければならない''
という条件の変密度・界面版である~\cite{Chorin1968,Guermond2006}．
面上の係数を
\[
  A_f =
  \begin{cases}
    \rho_q^{-1}, & f\subset\Omega_q\ \text{かつ同相面},\\
    0, & f\ \text{が相を横切る切断面},
  \end{cases}
\]
と書くと，投影恒等式は
\begin{equation}
  L_h^\mathrm{proj} = D_f A_f G_f,\qquad
  A_h^\mathrm{solve}=-L_h^\mathrm{proj},\qquad
  \bu_f^{n+1} = \bu_f^* - \Delta t\,A_fG_f(\delta p),\qquad
  D_f\bu_f^{n+1}=0
  \label{eq:same_face_projection_closure}
\end{equation}
である．ここで重要なのは，$A_f$，$G_f$，$D_f$ が
\textbf{PPE・速度補正・浮力残差の三者で同一}であることである．
式~\eqref{eq:buoyancy_gradient_residual_split} の
$\Phi_g\bnabla\rho'$ を面残差として評価する場合も，
節点で作った残差を事後的に面補間してはならない．
必ず同じ $G_f$ で hydrostatic ポテンシャルを差分し，同じ $A_f$ で圧力フラックス側へ戻す．
相を横切る面で $A_f=0$ とする分相係数を PPE だけに使い，
Corrector や浮力残差で混合密度係数を使うと，
投影形 $D_fA_fG_f$ で打ち消されるべき界面支持の発散が残る．
したがって「同じメトリック」だけでは不十分であり，
\textbf{同じ面係数・同じ面勾配・同じ面発散}まで含めて投影閉包と呼ぶ．

% ─────────────────────────────────────────────────────────────────────────────
\subsubsection{GFM ゴーストジェットによる位相ステッチ}
\label{sec:gfm_ghost_jet_stitch}
% ─────────────────────────────────────────────────────────────────────────────

界面 $\Gamma$ を跨ぐコンパクトステンシルを\textbf{単相化}するため，
production GFM では
ジャンプ条件 $[p]_\Gamma=J_p^\chi$，
$[\rho^{-1}\partial_n p]_\Gamma=0$
から\textbf{ゴースト面ジェット}
\begin{equation}
  \bigl(p_q^\ast,\,p_q'^\ast,\,p_q''^\ast\bigr)
    = \text{GFM extrapolate of phase }\bar q
      \text{ into } \Omega_q
  \label{eq:gfm_ghost_jet}
\end{equation}
を構築する（$\bar q$ は対相）．面 $f\subset\Gamma$ 付近のステンシルは
対相の値を\textbf{ゴーストジェットで置換}した上で FCCD を適用：
\begin{equation}
  (G_h^\FCCD p)_f^\text{stitched}
    = G_h^\FCCD\bigl(\,\{p_{q,i}\}_{i\in\Omega_q}
      \cup \{p_q^\ast(x_i)\}_{i\in\Omega_{\bar q}}\bigr).
  \label{eq:fccd_stitched_face}
\end{equation}
これにより\textbf{代数的な単相ステンシル}を各相で保ちつつ，
物理的不連続をジャンプクロージャ経由で自動的に注入する．

\noindent\textbf{ジャンプジェット構成}：
$\phi$（符号付距離，\S\ref{sec:ridge_eikonal}）と面位置 $x_f$ から
\begin{align}
  p_q^\ast(x_f)
    &= p_{\bar q}(x_f) - J_p^\chi\,\mathrm{sign}(\phi_f),
      \label{eq:gfm_value_jump}\\
  p_q'^\ast(x_f)
    &= \frac{\rho_q}{\rho_{\bar q}}\,p'_{\bar q}(x_f),
      \label{eq:gfm_gradient_jump}\\
  p_q''^\ast(x_f)
    &= p''_{\bar q}(x_f)
       - \bnabla_f\!\cdot(J_p^\chi\,\hat{\bm n}_{lg})\,\mathrm{sign}(\phi_f).
      \label{eq:gfm_curvature_jump}
\end{align}
これらのゴーストジェットは FCCD ステンシルを単相化するための片側延長であり，
affine jump の既知面ソースを別経路で二重に加えるものではない．
実装上は $j_{gl}/H_f$ 型の GFM 補正を
\S\ref{sec:split_ppe_pressure_jump_formulation} の $B_f(J_p^\chi)$ と同じ
face law として扱い，PPE RHS と corrector が共有する唯一のジャンプ寄与にする．
式~\eqref{eq:gfm_gradient_jump} は $J_n=0$ の production 形である．
一般の $J_n\ne0$ を許すなら，
$p_q'^\ast=(\rho_q/\rho_{\bar q})p'_{\bar q}+\rho_q J_n^{q\leftarrow\bar q}$
のように，向き付き flux-jump をゴースト導関数にも入れる必要がある．
本稿では非相変化・保存的 predictor を前提に $J_n=0$ へ固定する．
2 次値ジャンプの主要項は Laplace--Beltrami 分解により
$J_{ss}=\partial_s^2 J_p^\chi$ に縮退する．

\noindent\textbf{接線二階ジャンプ $J_{ss}$ の導出（界面座標 1 段階展開）：}
\textbf{前提仮定}（明示化）：
(H1) 滑らか延長 — 各相 RHS $\rho_q q$ は HFE（\S\ref{sec:hfe}）により
界面近傍で $C^2$ 級に延長されている；
(H2) flux-jump 消失 — $J_n\equiv[\rho^{-1}\partial_n p]_\Gamma=0$
（\S\ref{sec:gfm_ghost_jet_stitch}）；
(H3) 単位法線整合 — $|\bnabla\phi|=1$（リッジ Eikonal SDF 上）．
界面 $\Gamma$ 上で局所直交座標 $(s,n)$ を取る（$s$：接線弧長，
$n$：法線距離 $\phi$ と整合）．
中間補題：(H2) と運動量方程式表面張力ゼロ（\S\ref{sec:split_ppe_pressure_jump_formulation}）から
界面圧力ジャンプ $[p]_\Gamma=J_p^\chi$ を Heaviside 越しに
\textbf{1 階}展開した結果が式~\eqref{eq:gfm_curvature_jump} 第 1 式の
$[p_n]_\Gamma=\bnabla_f\!\cdot(J_p^\chi\hat{\bm n}_{lg})$%
\footnote{前提仮定 (H2) $[\rho^{-1}\partial_n p]_\Gamma=0$ は flux jump の連続性（GFM 求解条件），$[p_n]_\Gamma=\bnabla_f\!\cdot(J_p^\chi\hat{\bm n}_{lg})$ は表面張力下の curvature 駆動圧力ジャンプ（界面ソース項）であり，両者は独立に扱う；詳細は \S\ref{sec:split_ppe_pressure_jump_formulation}．}．
さらに $n$ 方向の Laplacian を分解すると：
\begin{equation}
  \nabla^2 p = \partial_n^2 p + \kappa_{lg}\,\partial_n p + \nabla_s^2 p,
  \label{eq:laplace_beltrami_normal_decomp}
\end{equation}
ここで $\nabla_s^2$ は界面 Laplace--Beltrami 演算子．
界面で $\partial_n p$ がジャンプを持つため，
$[\nabla^2 p]_\Gamma = [\partial_n^2 p]_\Gamma + \kappa_{lg}\,[\partial_n p]_\Gamma + [\nabla_s^2 p]_\Gamma$．
仮定 (H1) より両相 PDE $\nabla^2 p = \rho q$（$q\equiv\bnabla\!\cdot\bu^*/\Delta t$ の各相延長）
が界面まで滑らかに到達し，
$[\nabla^2 p]_\Gamma = [\rho q]_\Gamma$ は与えられた量．
$[\nabla_s^2 p]_\Gamma=\partial_s^2[p]_\Gamma=\partial_s^2 J_p^\chi$
から接線 2 階差分を取り，
1 階差分 $\kappa_{lg}[\partial_n p]_\Gamma=\kappa_{lg}\bnabla_f\!\cdot(J_p^\chi\hat{\bm n}_{lg})$ を
代入して整理すると，
\begin{equation}
  [\partial_n^2 p]_\Gamma
  = [\rho q]_\Gamma
    - \kappa_{lg}\,\bnabla_f\!\cdot(J_p^\chi\hat{\bm n}_{lg})
    - \partial_s^2 J_p^\chi.
  \label{eq:kappa_s_derivation}
\end{equation}
分相 PPE の片側で $[\rho q]_\Gamma$ を吸収する RHS 整合化と
仮定 (H3) $|\bnabla\phi|=1$ を用いれば，
2 階ジャンプ成分の主要項は接線 2 階微分
\begin{equation}
  J_{ss} \equiv \partial_s^2 J_p^\chi
  \label{eq:kappa_s_def}
\end{equation}
に縮退する（高次項 $\Ord{(h\kappa)^2}$ を含む詳細は
Liu--Fedkiw--Kang~\cite{Fedkiw1999} Appendix A の Laplace--Beltrami 分解と整合）．
本稿はこの $J_{ss}$ を分相 PPE の界面 BC 補正項として用いる．

% ─────────────────────────────────────────────────────────────────────────────
\subsubsection{PPE 整合性条件と affine 大域ゲージ}
\label{sec:ppe_phase_gauge_pin}
% ─────────────────────────────────────────────────────────────────────────────

相ごとに完全に切断された Neumann ブロック
$L_qp_q=S_q$ を解く比較用 formulation では，
各相で\textbf{可解条件}
\begin{equation}
  \sum_{i\in\Omega_q}\,S_{q,i}\,\Delta V_i = 0,
  \qquad S_{q,i} = \frac{1}{\Delta t}\,(D_h^\FCCD \mathbf{u}^*)_i
  \label{eq:split_ppe_compatibility}
\end{equation}
を満たす必要がある．しかし本稿の production 経路である
\texttt{affine\_jump} は，この「完全切断された二つの楕円問題」ではない．
界面を横切る face $f$ に対して，既知ジャンプを含む面 flux を
\begin{equation}
  F_{\Gamma,f}(p;j_{gl})
  =
  \alpha_f\{G_f(p)-B_f(j_{gl})\}
  \label{eq:split_ppe_affine_face_flux}
\end{equation}
として定義する cut-face 接続型の affine 演算子である．
ここで $B_f(j_{gl})=s_fj_{gl,f}/H_f$ は式~\eqref{eq:split_ppe_pressure_composition}
と同じ既知ジャンプ勾配であり，$H_f$ は面を挟む 2 点間の物理距離である．
この定義の核心は局所不変条件
\[
  p_\mathrm{high}-p_\mathrm{low}=s_f j_{gl,f}
  \quad\Longrightarrow\quad
  F_{\Gamma,f}(p;j_{gl})=0
\]
である．すなわち，すでに Young--Laplace ジャンプを満たす圧力場には
界面 face から人工 flux を発生させない．

\noindent\textbf{なぜ相別平均ゲージを課さないか：}
相別平均ゲージは，cross-phase face の圧力係数をゼロにした
切断ブロックに対して，各相の定数零モードを除去するための射影である．
これを affine face law にそのまま重ねると，式~\eqref{eq:split_ppe_affine_face_flux}
の同一係数性が壊れる．一般に
\[
  F^{\mathrm{split}}_{\Gamma,f}
  =
  \alpha^p_fG_f(p)-\alpha^J_fB_f(j_{gl})
\]
と regular 圧力側と既知ジャンプ側で係数を分けると，
切断ブロックでは $\alpha^p_f=0$ である．
$\alpha^J_f>0$ を残せば，ジャンプを満たす圧力でも
$F^{\mathrm{split}}_{\Gamma,f}=-\alpha^J_fB_f(j_{gl})\neq0$ となり，
affine の核心条件が破れる．一方 $\alpha^J_f=0$ まで落とせば，
界面ジャンプ駆動そのものが PPE から消え，旧問題
（capillary drive の相殺）に戻る．したがって
\texttt{affine\_jump} の production 演算子では，cut-face を
$G_f(p)$ と $B_f(j_{gl})$ の\textbf{同じ非零 face 係数}で扱い，
その連結グラフ上の零モードは全体定数 1 個だけになる．

\noindent\textbf{ゲージ条件：}
pressure jump は圧力の絶対値ではなく face-gradient 契約であり，
$p\mapsto p+C$ に対して $G_f(p)$ も $B_f(j_{gl})$ も変化しない．
したがって \texttt{affine\_jump} で取り除くべき零空間は
全体定数モードのみであり，実装では固定した大域ゲージを用いる．
相ごとの平均零拘束を追加すると，零空間でない
液相--気相間の相対定数差まで拘束してしまい，
Young--Laplace ジャンプとして保持すべき自由度を
数値ゲージで上書きすることになる．

\noindent\textbf{注意}：
圧力ジャンプ形式（\S\ref{sec:split_ppe_pressure_jump_formulation}）
で未知量が $\tilde p$（滑らか部分）となる場合，
反復解法のウォームスタートには$\tilde p$のみを渡す．
\textbf{regular 場として合成したジャンプ圧力}
をウォームスタートに用いると，不連続成分が次段 PPE に吸収され，
物理的な界面駆動が消える．第~\ref{sec:val_capillary}節の
root-cause 診断では，旧 \texttt{jump\_decomposition} 経路が
$p_\mathrm{base}\simeq -J$，$p_\mathrm{total}\simeq 0$ を作り，
$\operatorname{ptp}(p_\mathrm{total})/\operatorname{ptp}(J)=9.02\times10^{-4}$
までジャンプを代数的に相殺した．

% ─────────────────────────────────────────────────────────────────────────────
\subsubsection{圧力ジャンプ形式：CSF 体積力を回避するクリーン分離}
\label{sec:split_ppe_pressure_jump_formulation}
% ─────────────────────────────────────────────────────────────────────────────

本形式では CSF 体積力を使用せず，向き付き Young--Laplace
ジャンプ
\begin{equation}
  j_{gl} \equiv p_g-p_l = -\sigma\kappa_{lg}
  \label{eq:split_ppe_oriented_jump}
\end{equation}
を PPE と corrector が共有する face-gradient 契約として課す．
運動量方程式の表面張力項は\textbf{ゼロ}：
\begin{equation}
  \rho\frac{D\mathbf{u}}{Dt}
    = -\bnabla p + \bnabla\cdot(2\mu\mathbf{D}) + \rho\mathbf{g},
  \qquad p_g-p_l = j_{gl}.
  \label{eq:split_ppe_momentum_no_csf}
\end{equation}
ここで $\psi=1$ が液相，$\psi=0$ が気相であり，
$\kappa_{lg}=\nabla_\Gamma\!\cdot\hat{\bm n}_{lg}$ は液相から気相へ向く曲率である．
液滴では $\kappa_{lg}>0$ なので $j_{gl}<0$（液相圧が高い）となり，
気泡では $\kappa_{lg}<0$ なので $j_{gl}>0$（気相圧が高い）となる．

以下では全圧力 solve の表記で $J_p^\chi=j_{gl}$ と書く．
圧力増分 solve では，式~\eqref{eq:split_ppe_unknown_jump_contract} に従い
同じ場所を $\delta j_{gl}$ に置き換える．
実装上の未知量は regular 圧力 $\tilde p$ だけではなく，
面勾配演算子そのものを affine 化する：
\begin{align}
  s_f &= I_g(\mathrm{high})-I_g(\mathrm{low}),\qquad
        I_g=\mathbf{1}_{\psi<1/2}, \\
  H_f &= |\bm{x}_{\mathrm{high}(f)}-\bm{x}_{\mathrm{low}(f)}|, \\
  B_f^{\mathrm{nu}}(J_p^\chi) &= \frac{s_f\,J_{p,f}^\chi}{H_f}, \\
  G_{\Gamma,f}^{\mathrm{nu}}(p;J_p^\chi)
    &= G_f^{\mathrm{nu}}(p)-B_f^{\mathrm{nu}}(J_p^\chi).
  \label{eq:split_ppe_pressure_composition}
\end{align}
一様格子では $H_f=h$ で旧表記に戻る．非一様格子では
第~\ref{sec:grid}章の座標変換メトリクスが $G_f^{\mathrm{nu}}p$ を正しく評価するが，
界面を跨ぐ既知不連続 $B_f$ は滑らかな圧力場の微分ではない．
したがって新しい非一様 PPE ソルバは不要だが，
面ごとの $H_f$，符号 $s_f$，同じ face 係数を持つ
affine face-jump 演算子として接続する必要がある．
したがって PPE の界面項は
\begin{equation}
  D_f^{\mathrm{nu}}\alpha_fG_{\Gamma,f}^{\mathrm{nu}}(p;J_p^\chi)
  =D_f^{\mathrm{nu}}\alpha_fG_f^{\mathrm{nu}}(p)
   -D_f^{\mathrm{nu}}\alpha_fB_f^{\mathrm{nu}}(J_p^\chi)
  \label{eq:split_ppe_affine_rhs}
\end{equation}
であるから，投影式を
$D_f^{\mathrm{nu}}\alpha_fG_{\Gamma,f}^{\mathrm{nu}}=q$ と書けば
\[
  D_f^{\mathrm{nu}}\alpha_fG_f^{\mathrm{nu}}(p)
  =
  q+D_f^{\mathrm{nu}}\!\left(\alpha_fB_f^{\mathrm{nu}}(J_p^\chi)\right)
\]
となり，RHS にだけ既知量を注入する．corrector では同じ面データを用いて
\begin{equation}
  u_f^{n+1}
  =
  u_f^\ast
  -\Delta t\,\alpha_f
  \left(G_f^{\mathrm{nu}}\chi-B_f^{\mathrm{nu}}(J_p^\chi)\right)
  \label{eq:split_ppe_affine_corrector}
\end{equation}
を適用する．この符号は，液相 low / 気相 high の
二セル界面で $p_\mathrm{high}-p_\mathrm{low}=J_p^\chi$ なら
$G_{\Gamma,f}=0$ となる局所不変条件で固定される．
非一様格子 + 壁面 BC で速度補正が
$\mathcal{G}^\text{adj}$ へ切り替わる場合も同じ原則を保ち，
\begin{equation}
  G_{\Gamma,f}^{\mathrm{adj}}
  =
  G_f^{\mathrm{adj}}-B_f^{\mathrm{adj}},
  \qquad
  B_f^{\mathrm{adj}}\ \text{は}\ G_f^{\mathrm{adj}}\ \text{と同じ面平均・同じ}\ H_f\ \text{で構成する}.
  \label{eq:split_ppe_affine_adj}
\end{equation}
PPE で使う $D_f^{\mathrm{nu}}(\alpha_fB_f)$ と corrector で差し引く
$B_f$ が面・係数・距離を共有しない場合，Balanced--Force 条件は構造的に失われる．
この affine jump 形式は，通常の regular 圧力場
\begin{equation}
  p_\mathrm{legacy}=\tilde p + j_{gl}(1-\psi)
  \label{eq:split_ppe_legacy_pressure_composition}
\end{equation}
でジャンプを表す旧分解とは異なり，楕円 solve が $J$ を
自由圧力として吸収できない．したがって本形式が解く課題は，
CSF 体積力の評価位置ずれだけではなく，鋭い Young--Laplace
ジャンプを regular 圧力の自由度として相殺してしまう代数的欠陥である．
一方で，曲率 cap，HFE 延長誤差，長時間幾何エネルギー安定性は
この affine 化だけでは解決されず，§\ref{sec:val_capillary} の
別検証ゲートに残る．PPE には
\texttt{InterfaceStressContext} として
\texttt{pressure\_jump\_gas\_minus\_liquid}$=j_{gl}$ を渡し，
raw な $\sigma\kappa$ は保存しない．この契約は毛管波専用ではなく，
静止液滴・上昇気泡・平坦 Rayleigh--Taylor 界面の全てに同じ式を用いる．

% ─────────────────────────────────────────────────────────────────────────────
\subsubsection{再メッシュ再投影と PPE コンテキスト防護}
\label{sec:split_ppe_regrid_guard}
% ─────────────────────────────────────────────────────────────────────────────

界面適合格子の動的再構築（\S\ref{sec:grid_ale}）では
速度再投影 PPE がメインステップ PPE の\textbf{前段}で走る．
前ステップの圧力ジャンプ $J_p^\chi$（典型的には $j_{gl}^n$）を継承すると
不連続合成圧力が再投影内部で微分され，人工運動エネルギーを生む．
したがって再投影 PPE には界面ジャンプコンテキスト（$\psi,\kappa_{lg},\sigma,j_{gl}$）を渡さず，
メインステップ PPE のみが $[p]_\Gamma$ を保持する分離設計を採る．
